% V1013 P2.1 draft S1 -- Figure candidate for Part II intro (branch 2).
% Purpose: close the "charge/discharge label <-> delithiation direction" correspondence
% for the two electrodes side by side, so a reader can see in one picture why the
% *label* (charge/discharge) and the *physical content* fed to sigma_d (delithiation=+1)
% coincide for graphite but invert for LCO. Content is a direct diagram of the confirmed
% convention (verify10 B-2, f=+sigma_d; sec:lco-map; sec:lco-peak (a) footnote) -- no new
% physics, no numbers invented. Standalone-compilable snippet (uses the same
% tikz/positioning/arrows.meta libraries already loaded by graphite_ica_ch1_v1.0.13.tex).
\begin{figure}[t]
\centering
\begin{tikzpicture}[
  node distance=6mm and 10mm,
  box/.style={draw,rounded corners=2pt,minimum width=34mm,minimum height=8mm,align=center,font=\scriptsize,fill=black!4},
  hd/.style={font=\scriptsize\bfseries},
  ar/.style={-{Latex[length=1.6mm]},thick},
  lab/.style={font=\scriptsize,align=center}]

% -------- Graphite anode column --------
\node[hd] (ghd) at (-3.2,3.0) {graphite anode (this doc's main thread)};
\node[box] (gdis) at (-3.2,1.9) {cell label: discharge\\ $\sigma_d=+1$};
\node[box,below=of gdis] (gdel) {physical content:\\ delithiation ($\xi:0\to1$)};
\node[box,below=of gdel] (gup) {$V$ rises};
\draw[ar] (gdis) -- (gdel);
\draw[ar] (gdel) -- (gup);
\node[lab,below=3mm of gup] {label = physics\\(coincide)};

\node[box] (gchg) at (-3.2,-1.9) {cell label: charge\\ $\sigma_d=-1$};
\node[box,above=of gchg] (glith) {physical content:\\ lithiation ($\xi:1\to0$)};
\draw[ar] (gchg) -- (glith);
\draw[ar] (glith) -- (gup.south) [in=180,out=0];

% -------- LCO cathode column --------
\node[hd] (lhd) at (3.6,3.0) {LCO cathode (Part II)};
\node[box] (lchg) at (3.6,1.9) {cell label: charge\\ $\sigma_d=+1$ (delith. slot)};
\node[box,below=of lchg] (ldel) {physical content:\\ delithiation ($\xi:0\to1$)};
\node[box,below=of ldel] (lup) {$V$ rises};
\draw[ar] (lchg) -- (ldel);
\draw[ar] (ldel) -- (lup);
\node[lab,below=3mm of lup] {label $\ne$ physics\\(inverted)};

\node[box] (ldis) at (3.6,-1.9) {cell label: discharge\\ $\sigma_d=-1$};
\node[box,above=of ldis] (llith) {physical content:\\ lithiation ($\xi:1\to0$)};
\draw[ar] (ldis) -- (llith);
\draw[ar] (llith) -- (lup.south) [in=180,out=0];

% -------- shared bottom note --------
\node[lab,align=center] at (0.2,-4.6) {model input slot always reads \emph{delithiation $=+1$}\\(3 direction-dependent loci: polarization \S\ref{sec:pol}, branch \S\ref{sec:hys}, tail \S\ref{sec:tail} --\\ equilibrium species itself is direction-symmetric, $\xi\leftrightarrow1-\xi$)};

\draw[densely dashed,gray] (0.2,3.6) -- (0.2,-5.3);
\end{tikzpicture}
\caption{충$\cdot$방전 셀 라벨과 탈리튬화 방향의 대응 — 흑연 음극(왼쪽)은 라벨(방전)과 물리(탈리튬화)가 겹치지만, LCO
양극(오른쪽)은 \emph{충전}이 탈리튬화라 겹치지 않는다. 모델에 먹이는 방향 인자 $\sigma_d$ 는 두 전극 모두 셀 라벨이
아니라 탈리튬화 여부로 정해진다(f=$+\sigma_d$, \texttt{V1012\_P43\_verify10.md} \S2 확정 판정) — 이 슬롯이 갈라놓는
것은 분극$\cdot$분기$\cdot$꼬리 세 작용처뿐이고, 평형 종 자체는 $\xi\leftrightarrow1-\xi$ 대칭이라 방향 불변이다.}
\label{fig:p2-directionmap}
\end{figure}
